% =============================================================================
%  Notas de clase — Sesión 03: Sistemas de Numeración
%  Programación Avanzada — Universidad Panamericana
%  Compilar: pdflatex sesion03_sistemas_numeracion.tex
% =============================================================================
\documentclass[12pt, oneside]{article}
\usepackage{progavanzada}   % estilo "for dummies" alumno

\begin{document}

\portada{Sesión 03}{Sistemas de Numeración}{2026}

\tableofcontents
\newpage

% =============================================================================
\section{¿Por qué contamos en base 10?}
% =============================================================================

La respuesta corta: porque tenemos \textbf{diez dedos}.

Contar en base 10 (o \textit{sistema decimal}) significa que usamos diez
símbolos distintos —~0, 1, 2, \ldots, 9~— y que cada posición en un número
vale \textbf{diez veces} la posición de su derecha.

\begin{ejemplo}[El valor posicional en base 10]
\[
  9{,}324_{10}
  = 9 \times 10^3 + 3 \times 10^2 + 2 \times 10^1 + 4 \times 10^0
  = 9{,}000 + 300 + 20 + 4
\]
\end{ejemplo}

La posición más a la derecha vale $\text{base}^0 = 1$, la siguiente
$\text{base}^1$, luego $\text{base}^2$, y así sucesivamente.
Este patrón funciona \textbf{igual en cualquier base}; solo cambia el número
de símbolos disponibles y el peso de cada posición.

\begin{recuerda}
  Un sistema de numeración \textbf{posicional} es aquel donde el valor de
  un dígito depende de la posición que ocupa, no solo del símbolo en sí.
  El sistema romano (\textsc{iv}, \textsc{ix}\ldots) \textit{no} es
  estrictamente posicional; el nuestro, sí.
\end{recuerda}

% =============================================================================
\section{Otros sistemas a lo largo de la historia}
% =============================================================================

Los humanos no siempre contamos de diez en diez.
Distintas civilizaciones eligieron bases diferentes según su anatomía,
su astronomía o su necesidad de dividir cosas en partes iguales.

% -----------------------------------------------------------------------------
\subsection{Los mayas: base 20}
% -----------------------------------------------------------------------------

Los mayas usaban un sistema \textbf{vigesimal} (base 20).
La razón también es corporal: en lugar de contar solo con los dedos de las
manos, contaban además los \textbf{dedos de los pies}. Veinte extremidades,
veinte símbolos.

Su sistema tenía tres símbolos básicos: un punto (1), una raya (5) y un
caracol (0). Con ellos construían valores hasta 19 y luego cambiaban de
posición, igual que nosotros al llegar a 9.

\begin{tecnico}[Matemáticas mayas]
  El sistema maya fue uno de los primeros en incluir el \textbf{cero} como
  valor posicional, probablemente desde el siglo \textsc{iv} d.C.,
  siglos antes de que el concepto llegara a Europa.
  Para conocer más: Ifrah, G. \textit{Historia Universal de las Cifras}.
  Espasa Calpe, 1994.
\end{tecnico}

% -----------------------------------------------------------------------------
\subsection{Los babilonios: base 60}
% -----------------------------------------------------------------------------

¿Por qué un minuto tiene 60 segundos? ¿Por qué una hora tiene 60 minutos?
¿Por qué un círculo tiene 360 grados?

Todo eso viene de los \textbf{sumerios y babilonios}, que hace más de
4\,000 años usaban un sistema \textbf{sexagesimal} (base 60).

La elección no fue arbitraria: \textbf{60 es extraordinariamente divisible}.
Sus divisores son 1, 2, 3, 4, 5, 6, 10, 12, 15, 20, 30 y 60.
Eso significa que $\frac{1}{3}$, $\frac{1}{4}$ y $\frac{1}{5}$ de una hora
resultan en números enteros de minutos, algo muy útil para repartir, medir
campos o calcular calendarios.

\begin{consejo}[¿Por qué sobrevivió la base 60?]
  Los griegos heredaron la astronomía babilónica y con ella sus unidades.
  Ptolomeo dividió el círculo en 360 partes (grados), cada grado en
  60 \textit{partes primeras} (minutos) y cada minuto en 60
  \textit{partes segundas} (segundos). Esos nombres llegan hasta hoy.\\
  \textit{Referencia:} Neugebauer, O. \textit{The Exact Sciences in Antiquity}.
  Dover, 1969.
\end{consejo}

% -----------------------------------------------------------------------------
\subsection{La propuesta de la base 12}
% -----------------------------------------------------------------------------

Durante la Revolución Francesa (finales del siglo \textsc{xviii}), mientras
se diseñaba el sistema métrico decimal, varios matemáticos argumentaron que
sería más práctico adoptar la \textbf{base 12} (duodecimal) en lugar de la
base 10.

El argumento: \textbf{12 tiene más divisores enteros que 10}.

\begin{center}
\begin{tabular}{lll}
  \toprule
  \textbf{Base} & \textbf{Divisores} & \textbf{Fracciones exactas} \\
  \midrule
  10 & 1, 2, 5, 10           & $\frac{1}{2}$, $\frac{1}{5}$ \\[4pt]
  12 & 1, 2, 3, 4, 6, 12     & $\frac{1}{2}$, $\frac{1}{3}$, $\frac{1}{4}$, $\frac{1}{6}$ \\
  \bottomrule
\end{tabular}
\end{center}

Con base 12, un tercio o un cuarto serían tan ``limpios'' como la mitad lo es
hoy. Hoy en día la \textit{Dozenal Society of Great Britain} y la
\textit{Dozenal Society of America} siguen promoviendo esta idea.

\begin{consejo}
  Curiosamente, ya usamos la base 12 en la vida cotidiana: 12 meses,
  12 horas en el reloj, una docena de huevos, una gruesa (144 = $12^2$).
  La base 12 nunca desapareció del todo.\\
  \textit{Referencia:} \url{https://www.dozenalsociety.org.uk}
\end{consejo}

% =============================================================================
\section{¿Por qué las computadoras usan base 2?}
% =============================================================================

Una computadora está construida con millones de pequeños interruptores
electrónicos llamados \textbf{transistores}. Cada transistor solo puede
estar en uno de dos estados: \textbf{encendido} o \textbf{apagado}.

Por eso la base natural de cualquier computadora es la \textbf{base 2}:
dos estados, dos símbolos (0 y 1), llamados \textbf{bits}
(\textit{binary digits}).

\begin{recuerda}
  \textbf{¿Por qué no base 10?}\ \ Si quisiéramos representar diez estados
  distintos con voltaje eléctrico (0 V, 1 V, 2 V, \ldots, 9 V), cualquier
  ruido o variación mínima podría confundir un 3 con un 4. Con solo dos
  estados (bajo / alto), el circuito es \textbf{mucho más robusto y rápido}
  de fabricar y operar.
\end{recuerda}

\begin{tecnico}[Del bit al byte]
  \begin{itemize}
    \item \textbf{1 bit} — un dígito binario: 0 ó 1.
    \item \textbf{4 bits} — un \textit{nibble}: puede representar 0–15.
    \item \textbf{8 bits} — un \textit{byte}: puede representar 0–255.
    \item \textbf{32 / 64 bits} — tamaños típicos de las palabras
          (\textit{words}) que maneja el procesador de tu computadora.
  \end{itemize}
\end{tecnico}

% =============================================================================
\section{Sistema binario (base 2)}
% =============================================================================

En binario solo existen los dígitos \textbf{0} y \textbf{1}.
Cada posición vale el \textbf{doble} de la posición a su derecha
(potencias de 2).

\begin{center}
\begin{tabular}{ccccccccc}
  \toprule
  \textbf{Posición} & 7 & 6 & 5 & 4 & 3 & 2 & 1 & 0 \\
  \textbf{Valor}    & $2^7$ & $2^6$ & $2^5$ & $2^4$ & $2^3$ & $2^2$ & $2^1$ & $2^0$ \\
                    & 128 & 64 & 32 & 16 & 8 & 4 & 2 & 1 \\
  \bottomrule
\end{tabular}
\end{center}

% -----------------------------------------------------------------------------
\subsection{Binario → Decimal}
% -----------------------------------------------------------------------------

Multiplica cada bit por su peso posicional y suma.

\begin{ejemplo}[Convertir $1001_2$ a decimal]
\[
  1001_2
  = 1 \times 2^3 + 0 \times 2^2 + 0 \times 2^1 + 1 \times 2^0
  = 8 + 0 + 0 + 1
  = 9_{10}
\]
\end{ejemplo}

\begin{ejemplo}[Convertir $1100\,1000_2$ a decimal]
\[
  1100\,1000_2
  = 2^7 + 2^6 + 2^3
  = 128 + 64 + 8
  = 200_{10}
\]
\textit{Truco:} ignora los ceros; solo suma los pesos donde aparece un 1.
\end{ejemplo}

% -----------------------------------------------------------------------------
\subsection{Decimal → Binario}
% -----------------------------------------------------------------------------

Divide repetidamente entre 2 y anota los restos.
Los restos, leídos \textbf{de abajo hacia arriba}, forman el número binario.

\begin{ejemplo}[Convertir $25_{10}$ a binario]
\begin{center}
\begin{tabular}{r@{\;}c@{\;}l@{\quad}l}
  25 & $\div$ & 2 = 12 & resto \textbf{1} \\
  12 & $\div$ & 2 = 6  & resto \textbf{0} \\
   6 & $\div$ & 2 = 3  & resto \textbf{0} \\
   3 & $\div$ & 2 = 1  & resto \textbf{1} \\
   1 & $\div$ & 2 = 0  & resto \textbf{1} \\
\end{tabular}
\end{center}
Leyendo los restos de abajo hacia arriba: $25_{10} = 11001_2$.

\textit{Verificación:}
$16 + 8 + 1 = 25$. \checkmark
\end{ejemplo}

\begin{ejemplo}[Convertir $1521_{10}$ a binario]
Las divisiones sucesivas dan: $1521_{10} = 10111110001_2$.

\textit{Verificación:}
$1024 + 256 + 128 + 64 + 32 + 16 + 1 = 1521$. \checkmark
\end{ejemplo}

% =============================================================================
\section{Base 8 — Octal}
% =============================================================================

El sistema octal usa ocho símbolos: \textbf{0, 1, 2, 3, 4, 5, 6, 7}.
Cada posición vale el \textbf{óctuple} de la posición a su derecha
(potencias de 8).

\begin{ejemplo}[Convertir $4702_8$ a decimal]
\[
  4702_8
  = 4 \times 8^3 + 7 \times 8^2 + 0 \times 8^1 + 2 \times 8^0
  = 2048 + 448 + 0 + 2
  = 2498_{10}
\]
\end{ejemplo}

\begin{ejemplo}[Convertir $2498_{10}$ a octal]
\begin{center}
\begin{tabular}{r@{\;}c@{\;}l@{\quad}l}
  2498 & $\div$ & 8 = 312 & resto \textbf{2} \\
   312 & $\div$ & 8 = 39  & resto \textbf{0} \\
    39 & $\div$ & 8 = 4   & resto \textbf{7} \\
     4 & $\div$ & 8 = 0   & resto \textbf{4} \\
\end{tabular}
\end{center}
$2498_{10} = 4702_8$. \checkmark
\end{ejemplo}

\subsection{Octal $\leftrightarrow$ Binario: la relación directa}

3 bits representan exactamente los valores 0–7.
Por eso convertir entre octal y binario es inmediato: sustituye cada dígito
octal por su representación de \textbf{3 bits}.

\begin{center}
\begin{tabular}{cc}
  \toprule
  \textbf{Octal} & \textbf{3 bits} \\
  \midrule
  0 & 000 \\ 1 & 001 \\ 2 & 010 \\ 3 & 011 \\
  4 & 100 \\ 5 & 101 \\ 6 & 110 \\ 7 & 111 \\
  \bottomrule
\end{tabular}
\end{center}

\begin{ejemplo}[Convertir $4702_8$ a binario]
\[
  4702_8
  \;\to\;
  \underbrace{100}_{4}\;
  \underbrace{111}_{7}\;
  \underbrace{000}_{0}\;
  \underbrace{010}_{2}
  \;=\; 100\,111\,000\,010_2
\]
\end{ejemplo}

% =============================================================================
\section{Base 16 — Hexadecimal}
% =============================================================================

El hexadecimal usa \textbf{dieciséis} símbolos.
Como solo tenemos diez cifras (0–9), se completa con letras: A=10, B=11,
C=12, D=13, E=14, F=15.

\begin{center}
\begin{tabular}{cccc}
  \toprule
  \textbf{Decimal} & \textbf{Hex} & \textbf{Decimal} & \textbf{Hex} \\
  \midrule
   0 & 0 &  8 & 8 \\
   1 & 1 &  9 & 9 \\
   2 & 2 & 10 & A \\
   3 & 3 & 11 & B \\
   4 & 4 & 12 & C \\
   5 & 5 & 13 & D \\
   6 & 6 & 14 & E \\
   7 & 7 & 15 & F \\
  \bottomrule
\end{tabular}
\end{center}

\subsection{¿Por qué existe el hexadecimal?}

Porque escribir números binarios largos es \textbf{tedioso y propenso a
errores}. Un byte son 8 bits; un entero de 32 bits son 32 ceros y unos.
El hexadecimal es un \textbf{atajo} para representar binario de forma compacta.

La clave: \textbf{4 bits = 1 dígito hexadecimal} exactamente (pues
$2^4 = 16$).

\begin{ejemplo}[Convertir $\mathtt{EA23}_{16}$ a binario]
\[
  \mathtt{EA23}_{16}
  \;\to\;
  \underbrace{1110}_{E}\;
  \underbrace{1010}_{A}\;
  \underbrace{0010}_{2}\;
  \underbrace{0011}_{3}
  \;=\; 1110\,1010\,0010\,0011_2
\]
Sustituye cada dígito hex por su grupo de \textbf{4 bits}. Directo.
\end{ejemplo}

\begin{ejemplo}[Binario → Hexadecimal]
\[
  0111\,0000\,1001\,0011\,1100\,1011_2
  \;\to\;
  \underbrace{0111}_{7}\;
  \underbrace{0000}_{0}\;
  \underbrace{1001}_{9}\;
  \underbrace{0011}_{3}\;
  \underbrace{1100}_{C}\;
  \underbrace{1011}_{B}
  \;=\; 7093\mathtt{CB}_{16}
\]
Agrupa de derecha a izquierda en bloques de 4 bits y reemplaza.
\end{ejemplo}

\begin{consejo}
  En programación, el hexadecimal se escribe con los prefijos
  \texttt{0x} (C, C++, Python) o \texttt{\#} (colores en CSS/HTML).
  Por ejemplo, \texttt{0xFF} es $255_{10}$, y el color blanco en HTML
  es \texttt{\#FFFFFF}.
\end{consejo}

% =============================================================================
\section{Resumen de conversiones}
% =============================================================================

\begin{center}
\begin{tabular}{lcl}
  \toprule
  \textbf{Conversión} & \textbf{Método} & \textbf{Truco rápido} \\
  \midrule
  Cualquier base → decimal & Suma pesos posicionales  & $\sum d_i \times b^i$ \\
  Decimal → cualquier base & Divisiones sucesivas por $b$ & Leer restos de abajo arriba \\
  Binario $\leftrightarrow$ octal           & Grupos de 3 bits            & Directo, sin pasar por decimal \\
  Binario $\leftrightarrow$ hexadecimal     & Grupos de 4 bits            & Directo, sin pasar por decimal \\
  \bottomrule
\end{tabular}
\end{center}

% =============================================================================
\section{Ejercicios}
% =============================================================================

\begin{enumerate}
  \item Convierte $13{,}333_{10}$ a \textbf{binario}.
  \item Convierte $1001\,0110\,0100\,1100_2$ a \textbf{decimal}.
  \item Convierte $1111\,0\,1010\,1100\,0\,1010\,1011_2$ a \textbf{hexadecimal}.
  \item Convierte $3\mathtt{F2A1}_{16}$ a \textbf{binario}.
  \item Convierte $31037562_8$ a \textbf{hexadecimal}.
  \item Convierte $12300\,2130_4$ a \textbf{hexadecimal}.
\end{enumerate}

\begin{advertencia}
  Para los ejercicios 5 y 6 lo más sencillo es pasar primero a binario y
  luego agrupar en bloques de 4. Evita pasar por decimal si puedes.
\end{advertencia}

\vspace{1em}
\noindent\textbf{Problema de concurso}

Carina está escribiendo todos los números de uno en uno en su calculadora,
empezando en 1. Pero a la calculadora se le rompieron las teclas del 2 y del 7,
por lo que cualquier número que lleve esos dígitos se lo salta.
¿Cuál es el milésimo número que escribirá Carina?

\begin{consejo}[Pista]
  Si Carina no puede usar los dígitos 2 ni 7, en realidad está contando en
  un sistema de \textbf{ocho símbolos} (0, 1, 3, 4, 5, 6, 8, 9).
  ¿Qué relación tiene eso con una base que ya conoces?
\end{consejo}

% =============================================================================
\section{Resumen}
% =============================================================================

\begin{recuerda}
\begin{itemize}
  \item Contamos en base 10 por razones anatómicas; otras civilizaciones
        eligieron bases distintas (20 mayas, 60 babilónicos).
  \item Las computadoras usan base 2 porque los transistores tienen
        dos estados: encendido/apagado.
  \item \textbf{Binario → decimal:} suma los pesos posicionales donde hay un~1.
  \item \textbf{Decimal → base $b$:} divide repetidamente entre $b$
        y lee los restos de abajo hacia arriba.
  \item \textbf{Binario $\leftrightarrow$ octal:} grupos de 3 bits. \quad
        \textbf{Binario $\leftrightarrow$ hex:} grupos de 4 bits.
  \item Hexadecimal existe para escribir binario de forma compacta;
        es la notación más usada en programación de bajo nivel.
\end{itemize}
\end{recuerda}

\end{document}
